\par\Needspace{9\baselineskip}\bigskip
\pdfbookmark[2]{Erste Rechnung, [2]}{pallas-III-fragment2}
\begin{center}[2.]\end{center}
\begingroup
\newcommand{\IIILeaderPoint}{.\kern1.1pt}
\newcommand{\IIILeaderUpperLeft}{\rlap{\kern-.05em\raisebox{.5ex}{.}}\IIILeaderPoint}
\newcommand{\IIILeaderUpperRight}{\rlap{\kern.12em\raisebox{.5ex}{.}}\IIILeaderPoint}
\noindent\hspace*{\parindent}[Zur Berechnung der heliocentrischen Pallasörter, d. h. der Grössen \(r\) und \(v\), wurde das Elementensystem zum Grunde gelegt, welches die bisher beobachteten sechs Oppositionen möglichst gut darstellt und welches in der Disquisitio de elementis ellipticis Palladis art. 14 abgeleitet wurde (Band VI S. 24); reducirt auf das mittlere Aquinoctium 1803 Juni 30 (die Präcession in 181 Tagen ist gleich \(24'',\!83\)) ist dieses System das folgende:
\begin{center}
\fontsize{10}{12}\selectfont
\setlength{\tabcolsep}{3pt}
\begin{tabular}{@{}p{.68\linewidth}r@{}}
Epoche der mittlern Länge für den Meridian von Paris \IIILeaderPoint\IIILeaderPoint\IIILeaderPoint\IIILeaderPoint & \(260^{0}\,19'\,30'',\!55\)\\[3pt]
Mittlere tägliche tropische Bewegung \IIILeaderPoint\IIILeaderPoint\IIILeaderPoint\IIILeaderPoint\IIILeaderPoint\IIILeaderPoint\IIILeaderPoint\IIILeaderPoint\IIILeaderPoint\IIILeaderPoint\IIILeaderPoint\IIILeaderPoint\IIILeaderPoint\IIILeaderPoint\IIILeaderPoint\IIILeaderPoint\IIILeaderPoint\IIILeaderPoint & \(770'',\!50100\)\\[3pt]
Mittlere tägliche siderische Bewegung \IIILeaderPoint\IIILeaderPoint\IIILeaderPoint\IIILeaderPoint\IIILeaderPoint\IIILeaderPoint\IIILeaderPoint\IIILeaderPoint\IIILeaderPoint\IIILeaderPoint\IIILeaderPoint\IIILeaderPoint\IIILeaderPoint\IIILeaderPoint\IIILeaderPoint\IIILeaderPoint\IIILeaderPoint & \(770'',\!36383\)\\[3pt]
Länge der Sonnennähe \IIILeaderPoint\IIILeaderPoint\IIILeaderPoint\IIILeaderPoint\IIILeaderPoint\IIILeaderPoint\IIILeaderPoint\IIILeaderPoint\IIILeaderPoint\IIILeaderPoint\IIILeaderPoint\IIILeaderPoint\IIILeaderPoint\IIILeaderPoint\IIILeaderPoint\IIILeaderPoint\IIILeaderPoint\IIILeaderPoint\IIILeaderPoint\IIILeaderPoint\IIILeaderPoint\IIILeaderPoint\IIILeaderPoint\IIILeaderPoint\IIILeaderPoint\IIILeaderPoint\IIILeaderPoint\IIILeaderPoint\IIILeaderPoint & \(121^{0}\,8'\,33'',\!37\)\\[3pt]
Länge des aufsteigenden Knotens \IIILeaderPoint\IIILeaderPoint\IIILeaderPoint\IIILeaderPoint\IIILeaderPoint\IIILeaderPoint\IIILeaderPoint\IIILeaderPoint\IIILeaderPoint\IIILeaderPoint\IIILeaderPoint\IIILeaderPoint\IIILeaderPoint\IIILeaderPoint\IIILeaderPoint\IIILeaderPoint\IIILeaderPoint\IIILeaderPoint\IIILeaderPoint\IIILeaderPoint\IIILeaderPoint & \(172\quad28\quad37,26\)\\[3pt]
Neigung der Bahn \IIILeaderPoint\IIILeaderUpperLeft\IIILeaderPoint\IIILeaderPoint\IIILeaderPoint\IIILeaderPoint\IIILeaderPoint\IIILeaderPoint\IIILeaderPoint\IIILeaderPoint\IIILeaderPoint\IIILeaderPoint\IIILeaderPoint\IIILeaderPoint\IIILeaderPoint\IIILeaderPoint\IIILeaderPoint\IIILeaderPoint\IIILeaderPoint\IIILeaderPoint\IIILeaderPoint\IIILeaderPoint\IIILeaderPoint\IIILeaderPoint\IIILeaderPoint\IIILeaderPoint\IIILeaderPoint\IIILeaderPoint\IIILeaderPoint\IIILeaderUpperRight\IIILeaderPoint\IIILeaderPoint & \(34\quad37\quad28,35\)\\[3pt]
Excentricität \(=\sin14^{0}\,9'\,59'',\!79\) \IIILeaderPoint\IIILeaderPoint\IIILeaderPoint\IIILeaderPoint\IIILeaderPoint\IIILeaderPoint\IIILeaderPoint\IIILeaderPoint\IIILeaderPoint\IIILeaderPoint\IIILeaderPoint\IIILeaderPoint\IIILeaderPoint\IIILeaderPoint\IIILeaderPoint\IIILeaderPoint\IIILeaderPoint\IIILeaderPoint\IIILeaderPoint\IIILeaderPoint\IIILeaderPoint\IIILeaderPoint\IIILeaderPoint & \(0,2447424\)\\[3pt]
Logarithm der halben grossen Axe \IIILeaderPoint\IIILeaderPoint\IIILeaderPoint\IIILeaderPoint\IIILeaderPoint\IIILeaderPoint\IIILeaderPoint\IIILeaderPoint\IIILeaderPoint\IIILeaderPoint\IIILeaderPoint\IIILeaderPoint\IIILeaderPoint\IIILeaderPoint\IIILeaderPoint\IIILeaderPoint\IIILeaderPoint\IIILeaderPoint\IIILeaderPoint\IIILeaderPoint & \(0,4222071.\)
\end{tabular}
\end{center}
\par
\SourceLeaf{481}
Die heliocentrischen Jupitersörter, zur Berechnung der Grössen \(r',\,v',\,\beta'\), kann man aus den \textsc{Bouvard}schen Tafeln entnehmen. Da indessen die Berechnung von 61 Örtern nach diesen Tafeln eine sehr grosse Mühe erfordern würde, und auch die allergrösste Schärfe in den Jupitersörtern nicht erfordert wird, so kann man sich des folgenden Kunstgriffs*)\HistoricalUnmarkedNote{[*) Vgl. \textsc{Gauss} n \textsc{Olbers} 26. Nov. 1810; S. 415 dieses Bandes.]} bedienen, der im Wesentlichen darin besteht, dass man sich für die erste und zweite Hälfte des in Betracht kommenden Zeitraums von 3000 Tagen je ein elliptisches Elementensystem verschafft, das die Bewegung Jupiters genau genug darstellt.
\par
Dazu entnimmt man aus den \textsc{Bouvard}schen Tafeln die heliocentrischen Örter des Jupiter für die Anfangsepoche 1803 Juni 30 (Paris) \(=1803+181^{d}\), für die Endepoche 1811 September 16 \(=1803+3181^{d}\) und für die in der Mitte liegende Epoche 1807 August 8 \(=1803+1681^{d}\), nemlich:
\begin{center}
\setlength{\tabcolsep}{6pt}
\begin{tabular}{r|c|c|r}
 & helioc. Länge & Breite & \(\log r'\)\\[3pt]
\(1803.\ 181^{d}\) & \(188^{0}\,32'\,23'',\!3\) & \(+1^{0}\,18'\,52'',\!6\) & \(0,7367964\)\\[3pt]
\(1681\) & \(307\quad56\quad19,2\) & \(-0\quad38\quad49,7\) & \(0,7061432\)\\[3pt]
\(3181\) & \(82\quad47\quad10,8\) & \(-0\quad21\quad19,8\) & \(0,7081999.\)
\end{tabular}
\end{center}
\par
Die Reduction auf die Ekliptik beträgt für die drei Örter resp. \(+0'',\!2\), \(-23'',\!4\), \(+14'',\!1\). Es ergeben sich also die den drei Epochen entsprechenden Längen in der Bahn, reducirt auf das mittlere Äquinoctium 1803 Juni 30, wenn man die Präcession in \(3000^{d}\) zu \(411'',\!5\) annimmt, resp. gleich \(188^{0}\,32'\,23'',\!1\), \(307^{0}\,53'\,16'',\!8\), \(82^{0}\,40'\,5'',\!2\).
\par
Aus den beiden Örtern für \(1803.181^{d}\) und \(1803.1681^{d}\) kann nun das Elementensystem I bestimmt werden, das zur Berechnung der Jupitersörter für die erste Hälfte des Zeitraums dient, und aus den beiden letztern Örtern findet sich das Elementensystem II für die zweite Hälfte unseres Zeitraums. Und zwar geben zunächst die heliocentrischen Längen und Breiten nach der in Art. 110 der Theoria motus auseinander gesetzten Methode
\[
\begin{array}{lll}
\text{System \rlap{\kern-.08em\raisebox{1.7ex}{.}}I.}&\AscendingNode'=98^{0}\,23'\,47''&i'=1\quad18\quad53\\[5pt]
\text{System II.}&\AscendingNode'=98^{0}\,22'\,16''&i'=1\quad18\quad49
\end{array}
\]
\par
Ferner ergeben die Längen in der Bahn und die Radien Vectoren nach der in der Theoria motus Artt. 88—96 beschriebenen Methode die übrigen Elemente:
\begin{center}
\fontsize{10}{12}\selectfont
\setlength{\tabcolsep}{3pt}
\begin{tabular}{@{}p{.39\linewidth}rr@{}}
 & System I & System II\\[3pt]
\raggedright Epoche der mittlern Länge 1807 August 8 \IIILeaderPoint\IIILeaderPoint\IIILeaderPoint\IIILeaderPoint\IIILeaderPoint & \(300^{0}\,45'\,35''\) & \(300^{0}\,38'\,14''\)\\[4pt]
\raggedright Mittlere tägliche siderische Bewegung \IIILeaderPoint\IIILeaderPoint\IIILeaderPoint\IIILeaderPoint\IIILeaderPoint\IIILeaderPoint\IIILeaderPoint\IIILeaderPoint & \(299'',\!1063\) & \(299'',\!1053\)\\[4pt]
\raggedright Länge der Sonnennähe \IIILeaderPoint\IIILeaderPoint\IIILeaderPoint\IIILeaderPoint\IIILeaderPoint\IIILeaderPoint\IIILeaderPoint\IIILeaderPoint\IIILeaderPoint\IIILeaderPoint\IIILeaderPoint\IIILeaderPoint\IIILeaderPoint\IIILeaderPoint\IIILeaderPoint\IIILeaderPoint\IIILeaderPoint\IIILeaderPoint\IIILeaderPoint\IIILeaderPoint & \(12^{0}\,3'\,14''\) & \(12^{0}\,12'\,17''\)\\[4pt]
\raggedright Logarithm der Excentricität \IIILeaderPoint\IIILeaderPoint\IIILeaderPoint\IIILeaderPoint\IIILeaderPoint\IIILeaderPoint\IIILeaderPoint\IIILeaderPoint\IIILeaderPoint\IIILeaderPoint\IIILeaderPoint\IIILeaderPoint\IIILeaderPoint\IIILeaderPoint\IIILeaderPoint\IIILeaderPoint & \(8,686027\) & \(8,687916\)\\[4pt]
 & \(e=\sin2^{0}\,46'\,54''\) & \(=\sin2^{0}\,47'\,38''\)\\[4pt]
\raggedright Logarithm der halben grossen Axe \IIILeaderPoint\IIILeaderPoint\IIILeaderPoint\IIILeaderPoint\IIILeaderUpperLeft\IIILeaderPoint\IIILeaderPoint\IIILeaderPoint\IIILeaderPoint\IIILeaderPoint\IIILeaderPoint & \(0,716256\) & \(0,716257.\)
\end{tabular}
\end{center}
\par
Offenbar muss der Jupitersort für \(1803.1681^{d}\), aus beiden Systemen berechnet, sich als ein und derselbe erweisen.]
\endgroup
